\section{Introduction}
\label{sec:intro}

Modern software development demands constant movement across tools: a task
management board to plan work, a chat platform to coordinate with teammates,
and a code-hosting service to review pull requests and track issues. Each
transition carries a context-switching cost---the cognitive overhead of
re-establishing mental state in a new tool---that fragments attention and
increases workload~\cite{mark2008cost,czerwinski2004diary}. This fragmentation
is compounded when information lives in separate silos: acting on a chat
message that references a board card requires the developer to open two
applications and mentally merge their contents. Ko et
al.~\cite{ko2007information} documented this cross-tool information-seeking
overhead in collocated software teams; LaToza et al.~\cite{latoza2006maintaining}
found it a primary driver of knowledge loss and coordination breakdowns.

Integrated developer workspaces---products that co-locate boards, chat, and
code review---reduce this overhead by design, and they also create a new
possibility: an embedded AI agent that can \emph{observe and act} across all
surfaces at once. How much of that workspace such an agent should be allowed to
see, and what that breadth buys the user, remains thin in the empirical record.
Prior AI-productivity research has focused on single-surface tasks
(code generation, text drafting), leaving the cross-surface coordination layer
largely unstudied~\cite{chen2021codex,noy2023}. Agentic frameworks such as
ReAct~\cite{yao2023react} and Reflexion~\cite{shinn2023reflexion} show strong
performance on synthetic benchmarks but have not been evaluated with human
participants on realistic multi-surface workflows.

A second open question concerns \emph{context breadth}: how much workspace state
should an agent be allowed to observe? Wider context may improve agent relevance
but raises concerns about auditability and trust. Existing systems gate context
via prompt instructions, which are opaque and unverifiable---there is no way to
audit, after the fact, exactly what information an agent received during a
specific turn.

\textbf{This paper makes three contributions:}
\begin{enumerate}
  \item \textbf{\Weld system.} We describe \Weld, an open-source all-in-one
    developer workspace (Kanban board, team chat, GitHub integration) with an
    embedded study mode that provisions isolated per-participant workspaces,
    implements balanced condition assignment, and logs fine-grained behavioral
    telemetry~(\S\ref{sec:system}).

  \item \textbf{Code-enforced context scoping (\ContextGate).} We introduce the
    \ContextGate, a server-side module that assembles agent context from a
    condition-specific database allowlist, with a per-turn audit snapshot
    (\texttt{agent\_context\_sent}) that proves---from the server log alone---
    exactly which information sources a given agent turn could
    access~(\S\ref{sec:system}).

  \item \textbf{Empirical study.} We report a hybrid within/between-subjects
    user study ($N$~=~\emph{TBD}) that, against each participant's own no-agent
    baseline, crosses three agent-context conditions (\CI\ chat-only advisory;
    \CII\ chat+board action; \CIII\ chat+board+GitHub full) with three task
    types (comprehension, single-surface, cross-surface) to ask whether broader
    context makes users more effective (RQ1), how it shifts their reliance on
    versus verification of the agent (RQ2), and whether their trust tracks the
    agent's measured effectiveness (RQ3)~(\S\ref{sec:study},
    \S\ref{sec:results}).
\end{enumerate}

The manipulated variable is the agent's context breadth alone. We make no claim
that an all-in-one workspace outperforms separate tools---this design has no
fragmented-tooling arm to test that---and treat consolidation as the apparatus
and motivation, not as a hypothesis.

The remainder of the paper is structured as follows.
Section~\ref{sec:related} surveys related work on tool fragmentation, AI
productivity agents, and context management.
Section~\ref{sec:system} describes \Weld and its study infrastructure.
Section~\ref{sec:study} presents the study design.
Section~\ref{sec:results} reports results.
Section~\ref{sec:discussion} discusses implications, and
Section~\ref{sec:conclusion} concludes.
